\begin{abstract}
Interactive latency and memory-pressure stalls on modern mobile and embedded platforms arise from a structural mismatch: applications run on a managed runtime that governs a Java heap through garbage collection, while large native (C/C++) components allocate on an unmanaged heap that the runtime cannot see, and the kernel reclaims memory page-by-page from fault addresses without any notion of application-level heap intent. We characterize this mismatch empirically on commercial devices and isolate two root causes: a \emph{visibility gap}, in which the Android Runtime (ART) is unaware of native heap growth until a physical page fault forces a reaction, and a \emph{temporal gap}, in which runtime garbage collection and kernel reclamation act reactively and at different moments, so intervention arrives after pressure has already built. Guided by these observations, we present the Heap Memory Safeguard (HMS), an OS-level framework that co-locates runtime heap telemetry with kernel-level reclamation under a single coordinated control loop. The core of HMS is a heap-utilization index $H_u$ that fuses resident footprint with allocation velocity, allowing the system to distinguish transient bursts from sustained growth and to trigger \emph{targeted} reclamation before pressure escalates---rather than uniformly reclaiming more aggressively. HMS requires no application changes and reuses standard runtime, memcg, and hardware-tagging interfaces. On two commercial smartphones across camera, gaming, on-device inference, and multitasking workloads, HMS improves heap stability, shortens garbage-collection pauses, and lowers reclaim latency relative to runtime-only and kernel-only baselines, and our analysis attributes these gains to cross-layer coordination rather than to increased reclaim aggressiveness alone.
\end{abstract}

%% ---- CCS concepts (regenerate with the official ACM CCS tool before final submission) ----
\begin{CCSXML}
<ccs2012>
   <concept>
       <concept_id>10011007.10011006.10011050.10011017</concept_id>
       <concept_desc>Software and its engineering~Garbage collection</concept_desc>
       <concept_significance>500</concept_significance>
   </concept>
   <concept>
       <concept_id>10011007.10011006.10011050.10011051</concept_id>
       <concept_desc>Software and its engineering~Allocation / deallocation strategies</concept_desc>
       <concept_significance>500</concept_significance>
   </concept>
   <concept>
       <concept_id>10011007.10011074.10011099</concept_id>
       <concept_desc>Software and its engineering~Operating systems</concept_desc>
       <concept_significance>300</concept_significance>
   </concept>
   <concept>
       <concept_id>10010583.10010588.10010589</concept_id>
       <concept_desc>Computer systems organization~Embedded software</concept_desc>
       <concept_significance>300</concept_significance>
   </concept>
</ccs2012>
\end{CCSXML}

\ccsdesc[500]{Software and its engineering~Garbage collection}
\ccsdesc[500]{Software and its engineering~Allocation / deallocation strategies}
\ccsdesc[300]{Software and its engineering~Operating systems}
\ccsdesc[300]{Computer systems organization~Embedded software}

\keywords{Heap memory, memory management, garbage collection, cross-layer
coordination, mobile operating systems, Android runtime, kernel reclamation,
on-device inference}
